% 本文件由 LaTeX 排版 Agent 根据有效 translation.json 自动生成。
% author=Gregory of Nyssa; work_id=2903; title=论圣神——驳马其顿派

\chapter{论圣神——驳马其顿派}

% structure: p0001 source=h1 target=omitted reason=page-title-already-rendered-by-parent

对于愚昧的言论，一概不予回答，或许有失尊严；撒罗满智慧的劝诫似乎指向这一点：\Quote{不要照愚昧人的愚妄话回答他。}但恐怕由于我们的沉默，错误会战胜真理，异端的腐疮会侵入其中，摧毁信仰纯正的言语。因此，我认为必须作答——不是照这些人的愚妄，他们向我们的宗教提出这样的反对，而是为纠正他们堕落的思想。因为上述\WorkTitle{箴言}中的劝诫，我想，不是要求我们沉默，而是要纠正那些显露愚妄行为的人；也就是说，我们的回答不应停留在他们愚蠢的观念层面，而是要推翻那些关于教义的无知和受骗的观点。\par

那么，他们对我们提出什么指控？他们指责我们亵渎神明，因为我们对于圣神持有崇高的概念。凡是我们遵循教父的教导所承认关于圣神的，他们按照自己的理解，拿来作为攻击我们的把柄，谴责我们亵渎。例如，我们承认圣神与圣父和圣子同尊，以至于在一切可思想、可命名、虔诚可归于神性本质的事物上，他们之间毫无差异。我们承认，除了在位格上以独特属性被默观外，圣神确实来自天主，并属于基督，正如圣经所说；但他既不像圣父那样无始，也不像圣子那样是独生子；在某些独特的性质上要分别看待，而在其他一切方面，如我方才所言，他与圣父圣子完全等同。但我们的对手断言，他与圣父圣子没有任何生命的共融；因本质的差异，他在每一点上都低等、较小：在能力、荣耀、尊严，总之在言语或思想上我们归于天主性的一切事上都是如此；因此，他们在荣耀中无分，也无权享有同等尊荣；至于能力，他只有足够用于分派给他的部分活动的能力；他与创造的力量完全无关。\par

这就是他们对圣神的观念；其逻辑结论是，圣神本身没有我们虔诚在言语或思想中归于神性本质的任何标志。那么，我们该如何论证？我们不回答任何新的东西，不回答我们自己发明的东西，尽管他们挑战我们这样做；我们依靠圣经关于圣神的见证，从中我们得知圣神是天主的，应被称为天主。如果他们承认这一点，且不反驳受默感的话语，那么他们虽然急于与我们争辩，也必须告诉我们，他们为什么与我们争辩，而不与圣经争辩。我们所说的与圣经所说的并无不同。——但在神性本质本身，一旦我们相信了它，我们看不出圣经教导或我们常识所提示的任何区分；即那种会按照强弱程度在神圣超越的本质内部划分，使之因或多或少而改变自身的区分。因为我们坚信它是单纯、一致、非复合的，我们看不出其中有任何异质部分复杂或组合在一起，因此，一旦我们的心智把握了神性的概念，我们就凭借这名称本身的含义，接受了其中一切可想象适合天主性的事物的完美。事实上，天主性在每一方面都展现出完美，只要善可被找到。如果它在任何一点上欠缺、不完美，那么神性的概念就会在那一点上受损，以至于它不能在那方面是或被称为神性；因为对于一个不完美、有缺陷、需要外部添加的事物，我们怎能用这个词呢？\par

我们可以用物质实例来证实我们的论点。火自然会使触摸者感到热，它的所有部分都是如此；一部分并不比另一部分热；但只要它是火，它就表现出与自身不变的同一，在完全相同的活动中绝对一致；如果在任何部分冷却了，那部分就不再能称为火；因为随着它的发热活动变为相反的，它的名称也改变了。水、空气和构成宇宙的每一元素也是如此；每种元素都有相同的规定，不容许有过多或不足的概念；例如，水不能被称为更多或更少的水；只要它保持同等的湿度，水的名称就适用；但一旦它变为相反的性质，它的名称也必须改变。空气那种轻盈、浮力大、\Quote{敏捷}的性质也在每一部分显现；而那些稠密、沉重、向下坠落的，则从\Quote{空气}这个词的含义中消失。同样，只要天主性在一切可能归于它的虔诚属性上保持完美，凭借这种在一切善上的完美，它就不辜负其名称；但如果任何有助于完美概念的东西被减去，那么在该点上天主性的名称就被歪曲，不再适用于该主体。我们不能把一个干燥的东西称为水，不能把性质冰冷的东西称为火，不能把僵硬坚硬的东西称为空气，也不能把一个不直接蕴含完美概念的东西称为神性；事实上，后者更不可能。\par

如果圣神确实——不只是名义上——被圣经和我们的教父们称为\OriginalTerm{神圣的}{Divine}，那么那些反对圣神光荣的人还有什么理由呢？祂是神圣的，是绝对美善、全能、智慧、光荣、永恒的；祂是这一切能够提升我们思想以达其存在之伟大的事物。这些属性的主体的单一性证明祂并非仅按度量拥有它们，仿佛我们可以想象祂在其本质中是另一物，而因上述品质的临在变成了另一物。那种情况是那些具有复合性质的存有所特有的；而圣神在每一方面都同样是单一而纯一的。这是众所公认的；否认这一点的人并不存在。那么，既然祂的存有只有一个单纯且单一的定义，祂所拥有的善并非后天获得的善；而是，无论祂可能还有别的什么，祂本身就是美善、智慧、能力、圣化、正义、永恒、不朽，以及一切崇高且超越其他名称的名称。那么，那些不惧怕针对亵渎圣神所宣告的可怕判词的人，竟主张如此一位存有不拥有光荣，这是怎样的心态呢？因为他们明确提出了那论断：我们不应相信祂应受光荣；虽然我不知道他们凭什么认为不承认本质上光荣者的真实本性是有益的。\par

因为他们在辩护中不能援引这样的理由：祂是由主在次序上第三位交给门徒的，因此祂远离了我们关于神性的理想。在每处行善的活动毫无减弱或变化的情况下，认为数字次序是任何减弱或本质变化的标志，这是多么不合理！就像一个人看到三支火炬上分别燃烧的火焰（我们假设第三支火焰是由第一支火焰传递到中间，然后点燃末端火炬），却声称第一支的热量超过了其他；次一等的火焰在程度上有了变化，趋向较小；而第三支火焰根本不能称为火，尽管它像火一样燃烧、发光，并且做火所做的事。但如果第三支火炬确实是火，尽管是由先前的火焰点燃的，那么这些人的哲学是什么？他们因圣神在神圣的口唇上被提及于圣父和圣子之后，就渎神地认为可以轻视圣神的尊威。当然，如果在我们对圣神本质的概念中有任何低于神圣理想之处，他们证明祂不拥有光荣是做对了；但如果祂尊威的高超在每一点上都被感知，为什么他们吝于承认祂的光荣呢？好比一个人描述某人是人，却认为继续来说他是有理性的、会死的，或任何可归于人的其他属性是不安全的，因而取消了他刚刚承认的东西；因为如果他没有理性，他就根本不是人；但如果后者被承认，那么对\Quote{人}中已包含的概念怎么能犹豫呢？同样，对圣神，如果当人称祂为神圣时说的是真理，那么当人定义祂为可敬的、光荣的、美善的、全能的时候，也没有说谎；因为所有这些概念都与神性的观念一同被承认。所以他们必须接受两种选择之一：要么根本不称祂为神圣，要么不减去任何可归于神性的概念。因此，我们最确实必须将这两个思想彼此包含：即神圣本质，以及对那神圣超越本质的恰当思想、虔诚直觉。\par

既然已断言——且正确地断言——圣神属于神圣本质，而在\Quote{神圣}这一个词中，正如我们所说，包含了所有伟大的概念，那么，承认神性的人就潜在地承认了其余的一切——光荣、全能、一切表明超越的东西。拒绝在圣神身上承认这一点确实是可怕的；可怕，是因为应用于祂的对立术语列表中与上述术语相应者具有不一致性。我的意思是，如果一个人不承认光荣，他就必须承认不光荣；如果一个人排除祂的能力，他就必须接受其反面。对于尊荣、美善和任何其他超越性也是如此，如果它们不被接受，就必须承认其反面。\par

但如果所有人必须回避这一点，因为它甚至超出了最令人反感的亵渎，那么虔诚的心就必须接受关于圣神更高贵的名称和概念，并必须论及祂我们已经列举的一切：祂拥有尊荣、能力、光荣、美善，以及一切激发虔诚的事物。也必须承认这些实在并非不完美地或以其光辉品质的某种界限附着于祂，而是与其名称相称地无限。祂不应被视为拥有一定程度的尊威，然后变得不同；而是祂总是如此。如果你在诸世代之前开始数算，或者你凝视来世，你将发现尊威、光荣或全能没有任何减弱，以至于祂能因增加而增大，或因减少而缩小。祂完全而彻底地完美，在任何方面都不容许缩小。按照他们的假设，祂在哪方面被缩小，哪方面就会暴露于趋向羞辱祂的思想的入侵。因为不绝对完美的东西必须在某一点上被怀疑具有相反的特性。但如果连想到这一点都是极端心智错乱的标志，那么最好承认我们的信德：祂在一切美善中的完美是完全无限的、无界限的，在任何方面都不减损。\par

若关于祂的道理若加以深入探究，也应同样研究关乎圣子和圣父的道理。难道你不承认在一者身上与在另一者身上，荣耀的圆满性相同吗？我想凡思考者都会赞同。那么，既然圣父的尊荣是圆满的，圣子的尊荣也是圆满的，并且他们也承认了圣神的尊荣的圆满，这些新理论家为什么强加给我们，不允许我们在圣神身上享有与圣父和圣子同等的尊荣呢？至于我们，我们遵循上述考虑，发现我们不能思想也不能说，那无需增加即可圆满的事物，与另一物相比，是较不尊贵的；因为当我们拥有某物，由于其无瑕的圆满，理智无法发现其增加的可能，我也看不到它如何能发现任何减少的可能。但这些人否认同等尊荣，实际上是在断言其相对缺乏；当他们进一步沿着同样的思路，通过比较而产生的缩减，他们转移了虔敬对圣神所形成的一切概念；他们不承认祂在良善、全能或任何这类属性上的圆满。但如果他们回避如此公开的亵渎，并承认祂在所有美善属性上的圆满，那么这些聪明的人必须告诉我们，一个完美的事物如何能比另一个完美事物更完美或较不完美；因为只要完美的定义适用于它，该事物就不能在完美方面接受更多或更少。\par

那么，既然他们同意圣神是完全完美的，并且也承认真正的虔敬需要圣父和圣子也在一切美善上完美，他们有什么理由可以在已经承认了圣父之后又将祂夺走？因为夺走与圣父\Quote{同等的尊荣}，就是明确的证据，证明他们不认为圣神有分于圣父的完美。至于这尊荣本身的概念应用于天主时，他们想要将圣神排除在外，他们是什么意思呢？他们是指那种人与人之间的尊荣，藉言语和姿态表达尊重，以优先权等类似做法表示自己的谦逊，这些在今日愚蠢的风气中以\Quote{尊荣}的名义保持下来。但这一切都取决于实施者的善意；如果我们想象他们选择不实施的情况，那么在人类中没有人单纯因本性具有任何优势，以至于他必然比其他人更受尊荣；因为所有人都被相同的自然比例同等标记。这一点很明显，不容置疑。例如，我们看见今天因其所担任的职务而被群众视为尊荣对象的人，明天自己却成为那些施予尊荣的人之一，因为职务已转给他人。那么，他们认为天主的尊荣也是如此，以至于只要我们愿意施予，那天主的尊荣就被保留，但当我们停止施予时，它也因我们的意志而终止吗？荒谬的想法，也是亵渎！天主独立于我们，不会在尊荣上增长；祂永远同一；祂不能进入更好或更坏的状态；因为祂没有更好，也不接纳更坏。\par

那么，你用什么方式尊崇天主呢？你怎样举扬至高者？你怎样将光荣归于那超越一切光荣者？你怎样赞美那不可理解者？如果\BibleQuote{万民好像水桶中的一滴}\BibleRef{依 40:15}，如依撒意亚所说，如果所有活人一齐发出和谐统一的赞美之声，这区区一滴的礼物对本质上已光荣的祂能增加什么呢？\BibleQuote{诸天宣示天主的荣耀}\BibleRef{咏 19:1}，然而它们仍被视为祂价值的贫弱宣告者；因为祂的尊威被高举，并非达到诸天，而是高于那些诸天，诸天本身包含在天主的一小部分之内，被比喻地称为祂的\Quote{一掌}。而人，这脆弱短命的受造物，如此恰当地被比作\Quote{草}，\Quote{今天还在}，明天就不在了，他能相信自己能适当地尊崇天主吗？这就像有人从麻絮中点燃一根细丝，幻想那火花能增加太阳耀眼的光芒。请问，如果你真的想尊崇圣神，你用什么言语尊崇祂呢？说祂绝对不朽，无转变、无变化，永远美丽，永远不依赖于他人加予，按自己的意愿在万物中实行一切，圣善、引导、正直、公义、言语真实，\BibleQuote{探究天主的深奥事}\BibleRef{格前 2:10}，\BibleQuote{由圣父所发}\BibleRef{若 15:26}，\BibleQuote{从圣子领受}\BibleRef{若 16:14}，以及所有类似的事，你用这些以及类似的词语到底给了祂什么？你是提到祂所拥有的，还是用祂所没有的来尊崇祂？如果你证明祂所没有的，你的话就毫无意义，归于虚无；因为称苦的东西为\Quote{甜}的人，自己说谎，却没有赞许那该受责备的。然而，如果你提到祂所拥有的，这样的品质是本质性的，无论人承认与否；宗徒说，纵然我们不信，祂仍是信德的对象。\par

那么，这些人灵性上的降低和扩展是什么意思？他们热心于圣父的尊荣，给予圣子同等的分享，但在圣神方面却缩减恩惠；既然已经证明天主的固有价值不依赖于我们的任何意志，而始终不可分割地内在于祂，他们这样做是什么意思？他们的心胸狭隘和忘恩负义，在这意见中暴露无遗，而圣神本质上是尊贵的、荣耀的、全能的，以及我们所能想到的一切崇高，尽管他们反对。\par

\Quote{是的，}他们中有一人回答说，\Quote{但圣经教导我们，圣父是造物主，同样也教导我们\InnerQuote{万物是藉着圣子}而造的；然而天主的话关于圣神从未说过这类事；那么，又怎能将圣神置于与那位通过创造显示如此伟大权能者同等的尊位呢？}\par

对此我们该如何回答呢？他们心中的想法纯属空谈，因为他们想像圣神并非始终与圣父和圣子同在，而是随着情况的变化，时而独自被默观，时而又紧密地与他们结合。因为如果天地及一切受造之物确实是通过圣子源于圣父，却与圣神无关，那么当圣父与圣子共同创造时，圣神在做什么呢？祂是否忙于其他工作，因此未参与宇宙的建造？但若是这样，在世界被造之时，他们又能指出圣神的什么特殊工作呢？的确，想像在圣父通过圣子所产生的创造之外还有另一个创造，乃是愚昧的妄念。好吧，假设祂根本没有参与，而是由于贪图安逸、不愿劳苦，所以远离了忙碌的创造工作？\par

愿仁慈的圣神亲自宽恕我们这无根据的假设！这些理论家的亵渎迫使我们必须步步紧跟，以至于无意中用他们自己的幻象污染了我们的论述。合乎一切虔敬的观点如下：我们不应认为圣父曾与圣子分离，也不应认为圣子与圣神分离。正如除非我们的思想通过圣子被提升，否则无法上升至圣父；同样，除非藉着圣神，也无法说耶稣是主。因此，圣父、圣子、圣神唯有在一个完美的\OriginalTerm{天主圣三}{Trinitas}中被认识，彼此处于最紧密的连续与结合中，在一切受造物之前，在万世之前，在我们所能想像的任何事物之前。圣父永远是圣父，圣子在祂内，圣神与圣子同在。既然这些位格彼此不可分离，那么这些人试图用时间上的区分来分割这种不可分割性，并自信地断言\Quote{唯有圣父，藉着圣子，创造了万物}，而圣神在这创造之时完全不在场或不工作，这是何等愚蠢！如果祂不在场，他们必须告诉我们祂在哪里；既然天主涵盖万物，他们能否想像圣神有一处独立的居所，以至在创造过程中祂能保持孤立？如果祂在场，为何不活动？是因为不能，还是不愿工作？祂是自愿不参与，还是因为某种强力迫使祂离开？如果祂故意选择不活动，那么祂也必拒绝以任何其他可能的方式工作；而那位宣称\Quote{祂随己意\BibleRef{格前 13:6}在众人内行一切事}者，按他们的说法就成了说谎者。反之，如果这位圣神有工作的冲动，但某种压倒性的控制阻碍了祂的计划，他们必须告诉我们这阻碍的原因。是由于祂被嫉妒参与这些工程的荣耀，为的是保证对成功工程的钦佩不延及第三人？还是因为对祂帮助的不信任，仿佛祂的合作会导致现时的恶果？这些聪明人无疑为我们提供了这两种假设之一的根据；否则，如果嫉妒与天主性无关，正如在无误者身上不可能想到失败一样，那么他们这种将圣神的能力与世界创造的功效隔离开的狭隘观点又有何意义呢？他们本应通过更高超的观念摒弃卑微的人类思维方式，并进行与所论对象之崇高更相称的估算。因为宇宙的天主并不需要帮助，所以并非\Quote{藉着圣子}创造宇宙；独生天主的能力也不欠缺，所以并非\Quote{藉着圣神}行一切事；而是权能之泉是圣父，圣父的权能是圣子，那权能的神是圣神；整个受造界，无论有形无形，都是那神圣权能的完成之工。既然在天主性任何事物的构成中都不能想到辛劳（因为意愿的时刻即成为现实，计划立刻成为现实），我们有理由将一切因创造而产生之本性称为一种意志的运动、一种计划的推动、一种权能的传递，始于圣父，通过圣子前进，并在圣神内完成。\par

这就是我们以平常方式所持的观点；我们摒弃了敌手所有这些精巧的诡辩，相信并宣认：在一切行为与思想中，无论今世或来世，无论时间中或永恒中，圣神都被视为与圣父、圣子结合，在意志、能力或任何其他对至高美善的虔敬概念中毫无欠缺；因此，除了次序与位格的区分外，不应认为有任何差异；但我们主张，虽然祂在单纯序列中被数算为第三，在传递次序中为第三，但在其他一切方面，我们承认祂们不可分离的结合：在性体、尊荣、天主性、荣耀、威能、全能以及一切虔敬信仰中，祂们都不可分离。\par

但关于侍奉与崇拜，以及他们如此精心计算并突出强调的其他事情，我们这样说：圣神远超我们以纯粹人性的目的所能为祂所做的一切；我们的崇拜远不及应得的尊荣；而人类习俗中视为可敬的任何事物，都在某种程度上低于圣神的尊严；因为那本质上不可测量的，超越了那些以如此微小、有限且微不足道的给予能力献上一切的人。\newline{}那么，我们对此对那些赞同敬畏圣神为神圣且具有神圣本性的概念的人说这些。但如果他们中有人拒绝这一陈述，以及这一包含在神性之名本身的观念，并且说那种流传于众人之中、损害圣神伟大的话，即祂不属于创造者，而属于受造物，认为应当视祂为受造的本性而非神圣的本性，我们对这样的命题回答说：我们完全不理解如何能将提出它的人算作基督徒之列。\newline{}因为正如不能将未成形的胚胎称为人，而只能称为潜在的人，假定它完成而得以出生为人，但只要它仍处于未成形状态，它就并非人；同样，我们的理性不能承认一个未能就全部奥迹接受我们宗教真确形式的人为基督徒。\newline{}我们能听到犹太人相信天主，也相信我们的天主：甚至我们的主在福音中提醒他们，他们不承认别的天主，只承认独生子的父，\Quote{你们说他是你们的天主。}\newline{}那么，我们是否因为犹太人同意崇拜我们所敬拜的天主而称他们为基督徒？我也知道摩尼教徒四处夸耀基督的名。因为他们崇敬我们屈膝敬拜的名，我们是否因此将他们算在基督徒之列？\newline{}同样，那既相信圣父又接受圣子，却置圣神的威严于不顾的人，已\Quote{背弃了信德，比不信的人更坏}，并且否认他所肩负的基督之名。\newline{}宗徒命令天主的人要\Quote{成全}。现在，仅就\Emph{一般}的人而言，成全必须在于人性每一方面的完满，即具有理性、思想与知识的能力、分享动物生命、直立姿态、笑的能力、宽阔的指甲；如果有人称某人为人，却无法在他身上提供上述人性标志的证据，那么他如此称呼就是无价值的尊称。\newline{}同样，基督徒的标志在于他相信圣父、圣子和圣神；按真理奥迹塑造之人的形式在于此。\newline{}但如果他的形式以别的方式安排，我将不承认任何缺乏形式的东西的存在；当圣神未被列入信仰时，标志模糊，本质形式丧失，我们完满人性的特征标志发生改变。\newline{}因为的确\BibleBook{}{训道篇}的话说真了；你的异端者不是活人，而是\Quote{骨头}，他说，\Quote{在孕妇的胎中}；因为一个不将傅油与受傅者一同思想的人，怎能说是相信受傅者呢？\Quote{祂}，\Person{伯多禄}说，\Quote{天主以圣神傅了祂。}\par

这些毁灭圣神荣耀的人，将祂贬到受造世界，必须告诉我们那傅油是何种事物的象征。难道它不是王权的象征吗？什么？他们不相信独生子就其本性是一位王吗？那些没有一次永远用犹太人的\Quote{帕子}\BibleRef{格后 3:14}蒙住心的人，不会否认祂就是。\newline{}那么，如果圣子就其本性是王，而傅油是祂王权的象征，那么你的理性会推论出什么结果呢？那就是：傅油并非与那王权相异之物，因而圣神在天主圣三中也不被列为奇怪或外来的东西。因为圣子是王，而祂那活的、实现的和位格化的王权存在于圣神内，圣神傅了独生子，从而使祂成为受傅者，成为一切存在之物的王。\newline{}那么，如果圣父是王，独生子是王，而圣神就是王权，则同一的王权定义必须贯穿这天主圣三，而\Quote{傅油}的概念传递了隐藏的含义，即圣子与圣神之间没有分离的间隔。\newline{}因为正如在身体表面与油液之间，无论是理性还是感知都无法察觉任何中介之物，圣神与圣子的结合是如此不可分离；结果是，凡是要因信触摸圣子的人，必须在触摸的行动中首先遇到油；祂没有任何部分缺少圣神。\newline{}因此，对圣子主权的信仰，在那些持有它的人中，是通过圣神产生的；那些因信接近圣子的人，四面都遇到圣神。那么，如果圣子本质上是一位王，而圣神是那傅了圣子的王权尊严，有什么本质上、与自己相比的这种王权的减损可以设想呢？\par

我们再次以这种方式来看。\OriginalTerm{王权}{Kingship}确凿无疑地体现在对臣民的统治上。现在，什么是臣服于这位\OriginalTerm{君王}{King}的呢？\OriginalTerm{圣言}{the Word}无疑包含世代及其中的一切；\Quote{你的王国，}它说，\Quote{是世代的王国，}并且，藉着世代，它意指在它们中间、在无限空间里受造的一切实体，无论有形无形；因为在它们内，万物都是由那些世代的创造者所造的。那么，如果王权必须总是与君王一起被思考，而臣民世界被承认为与统治者世界不同，这些人如此自相矛盾，将\OriginalTerm{傅油}{Unction}当作表达那本性为君王者价值的术语，却贬低那傅油本身到臣民的等级，仿佛它缺乏这种价值，这是何等荒谬！如果它因本性而是一个臣民，那么为什么它被立为王权的傅油，并与\OriginalTerm{独生子}{Only-begotten}的君王尊威联系在一起？另一方面，如果统治的能力通过它被包含在王权的威严中而显示出来，那么有什么必要将一切都拖入平民和奴仆的低下境况，并与\OriginalTerm{受造的}{Created}臣民同列呢？当我们论及圣神而肯定两种状态时，我们不可能在两种情况下都说真话：\Emph{i.e.}，祂是统治的，以及祂是臣服的。如果祂统治，祂就不在任何主之下；但如果祂是臣服的，那么祂就不能与那为君王者同被理解。人被人认出是在人当中，天使在天使当中，万物各从其类；因此\OriginalTerm{圣神}{Holy Spirit}必须被相信只属于两个世界中的一个：统治的世界或低下的世界；因为在两者之间，我们的理性不能认识任何东西；不能想到任何在\OriginalTerm{受造的}{Created}和\OriginalTerm{非受造的}{Uncreated}边界上的自然属性的新发明，以至于既参与两者又不完全是两者；我们不能想象这样一种对立物的融合和焊接，即由受造和非受造的东西混合，两个对立物如此合并成一位，在这种情况下，那种奇怪混合的结果不仅是一个复合体，而且由不同且时间上不一致的元素组成；因为那从受造中取得位格者，确实后于那无受造而自存者。\par

如果，那么，他们宣称\OriginalTerm{圣神}{Holy Ghost}是两者混合的，他们必然会把那种混合视为先与后的事物；并且，按照他们，祂会先于自己；相反地，后于自己；从\OriginalTerm{非受造的}{Uncreated}祂会得到先性，从\OriginalTerm{受造的}{Created}得到后性。但是，在事物的本性中，这是不可能的；因此，最肯定真实的是对\OriginalTerm{圣神}{Holy Spirit}肯定这两个替代中只有一个，那就是非受造的属性；因为注意另一个替代所涉及的荒谬数量；我们在实际受造物中能想到的一切事物，由于它们都因一个创造行为而获得存在，在所有情况下都完全相等的等级和价值，那么有什么理由将圣神与受造物的其余部分分开，并将祂与圣父和圣子同列呢？那么，逻辑会关于祂发现这一点：那被视为非受造部分者，不是因创造而存在；或者，如果它是，那么它就没有比其同类受造物更大的能力，它不能与那\OriginalTerm{超越本性}{Transcendent Nature}相联合；另一方面，如果他们宣称祂是一个受造者，同时又具有超越受造物的能力，那么受造物将被发现与自己不一致，分为统治者和被统治者，以致部分为施益者，部分为受益者，部分为圣化者，部分为被圣化者；而那一整套我们相信由圣神为受造物提供的祝福，都存在于祂内，丰涌而出，倾注于他人，而受造物仍需要从那源头涌出的帮助和恩宠，并仅仅作为施舍接受那些可以从同类受造物传递而来的祝福！那会是偏袒和以貌取人；我们知道在事物本性中没有这种不公，如同那些在实体方面毫无区别的存在者不应具有同等的权力；并且我认为，凡思考者都不会接受这种观点。要么祂不将任何事物赋予他人，如果祂本质上不拥有任何事物；要么，如果我们确实相信祂给予，就必须承认祂预先拥有那恩赐；这种赐予祝福而自身不需要如此外来帮助的能力，是\OriginalTerm{天主性}{Deity}独特而精妙的特权，而非其他。\par

那么，我们也来看这一点。在圣洗圣事中，我们由此获得的是什么？难道不是分享一个不再受死亡束缚的生命吗？我想，任何可以算作基督徒的人都不会否认这一点。那么，是什么？给予生命的权能是在那用来传递圣洗恩宠的水本身中吗？还是人人都清楚，这元素只被用作外在职务的工具，其本身对圣化毫无贡献，除非它先被圣化本身所转化？而赋予受洗者生命的乃是圣神，正如主亲自用祂的口论及祂所说的：\Quote{那赐予生命的是圣神}；但为使这恩宠得以完成，唯独凭信德所接受的祂并不赐予生命，而必须先相信我们的主，以便这活泼的恩赐临到信徒身上，正如主所说的：\Quote{祂愿意赐给谁生命，就赐给谁}。然而更进一步，由于这经由圣子施予的恩宠依赖于万有的\OriginalTerm{未生之源}{Ungenerate Source}，圣经因此教导我们，要先相信那生万有的圣父；这样，这赋予生命的恩宠，应当从这根源（如同从涌流生命的泉源）开始，通过那真实生命的\OriginalTerm{独生子}{Only-begotten}，借着圣神的运作，在适于领受的人身上得以完成。那么，如果生命在圣洗中来临，而圣洗在父、子、圣神的名内得以完成，这些人把这位生命的施予者视若无物，是什么意思呢？如果这恩赐是微不足道的，他们必须告诉我们比这生命更贵重的是什么。但若一切贵重之物都次于这生命，我是指那更高贵的生命，在动物身上没有份的生命，他们怎敢贬低如此巨大的恩惠，甚或贬低实际施予恩惠的那一位，并凭借他们的观念将祂贬低到受造世界，将祂与更高的神性世界分离呢？最后，如果他们坚持认为这生命的赐予是一件小事，并且在施予者身上没有什么伟大和可敬畏之处，他们为何不推断出这个观点本身必然导致的结论，即：我们必须认为，就连独生子和圣父本身，在祂们的生命中也并无伟大之处，而这生命与我们通过圣神所拥有的生命相同，由父通过子赐予我们？\par

因此，如果这些蔑视并攻击他们自己生命的人认为这恩赐是件小事，并据此决定轻视施予恩赐的那一位，那么让他们明白：他们不能只将自己的忘恩限制在一个位格上，而是必须将其不敬延伸到圣神之外的整个天主圣三。因为正如恩宠从父那里，通过圣子和圣神，源源不断地流到配得的人身上，同样，这种不敬也向后回流，从圣子传到万有的天主，从一个传到另一个。如果当一个人被轻视时，派遣他的那位也被轻视（然而人与派遣者之间的距离何其大！），那么那些如此侮辱圣神的人是何等有罪！也许这就是对\OriginalTerm{立法者的亵渎}{the blasphemy against our Law-giver}，为此定下了无法赦免的判决；因为在这其中，整个有福而神圣的存有也被侮辱了。正如虔诚敬拜圣神的人在祂身上看到独生子的光荣，并在这景象中看到无限天主的肖像，并通过那肖像在自己的认知中勾勒出原型的轮廓，同样，这个\OriginalTerm{蔑视圣神者}{contemner of the Spirit}，每当他提出任何胆大妄言反对圣神的光荣时，通过相同的推理，也将他的不敬延伸到圣子，并通过圣子延伸到圣父。因此，思想的人必须害怕，以免他们做出鲁莽之举，其结果将是行动者的彻底消灭；而他们在称呼中高举圣神时，甚至会在称呼之前就在思想中高举祂，因为言语不可能与思想齐头并进；然而，当一个人达到人类能力的最高极限，即心智所能达到的最高高度和宏伟观念时，即使那时也必须相信它远低于属于祂的光荣，符合圣咏中的话：\BibleQuote{在你们高举主我们的天主之后，你们还勉强朝拜祂脚下的脚凳}；而这尊贵之所以如此不可理解，别的原因，正是因为祂是圣的。\par

那么，如果人类能力的每一高度都低于圣神的伟大（因为这就是圣言在\Quote{脚凳}比喻中所指的意思），那么那些认为自身有如此巨大的能力，以至于可以定义一位无价之宝的价值的人是多么虚荣！因此他们宣称圣神不配一些与价值观念相关的东西，好像他们自己的能力远超过按他们估计的圣神的能力。多么可悲、多么可怜的疯狂！他们谈论这些话时，不明白自己是什么，也不明白他们所傲慢对抗的圣神是什么。谁将告诉这些人，人是\BibleQuote{去而不返的魂}\BibleRef{智 16:14}，在母腹中由不洁的受孕所建造，全部归于不洁的尘土；继承一个如同草芥的生命；在生命的幻象中短暂绽放，然后就枯萎，一切花朵凋落消失；他们自己也不确定出生前是什么，以及将变成什么，他们的灵魂在肉体中停留时不知道自己的特殊命运？人就是如此。\par

反之，圣神首先因本质圣洁的属性而正是圣父——本质圣洁者——所是的；并且正如\OriginalTerm{独生子}{Only-begotten}之所是，圣神亦然；其次，祂因赋予生命、不朽、不变、永恒、正义、智慧、正直、主权、良善、能力、赐予一切美善之事——尤其生命本身——的能力而如此；又因无所不在、临于万物、充满大地、居于诸天、倾注于超自然能力、按各人之所应得充满万有、自身保持圆满、与所有堪当者同在、却与天主圣三不离不弃。祂永远\Quote{洞悉天主深奥的事}，永远从子\Quote{领受}，永远被\Quote{派遣}，却不相分离，且受\Quote{光荣}，而祂一直拥有光荣。事实上，显然，将光荣给予他人者，自身必拥有丰盛的光荣；因为一个缺乏光荣者怎能光荣他人？除非事物本身是光，它怎能显耀光的恩赐？因此，光荣的能力永不能由自身不是光荣、尊威、威严和伟大者所显耀。如今圣神确实光荣父和子。那位说\Quote{光荣我的，我必光荣他}\BibleRef{撒上 2:30}者并不说谎；主向父说\Quote{我已光荣了你}；祂又说\Quote{求你以我在创世以前同你所有的光荣光荣我}。神圣的声音回答说\Quote{我光荣了，还要光荣}。你看见光荣的循环圈从相似者移向相似者。子由圣神光荣；父由子光荣；子又从父得光荣；如此，\OriginalTerm{独生子}{Only-begotten}成为圣神的光荣。因为父将由什么得光荣，若不是子的真实光荣？子又将由什么得光荣，若不是圣神的威严？同样，信德完成这循环，并藉着圣神光荣子，藉着子光荣父。\par

若圣神的伟大如此，且一切道德美善、一切良善，既由天主藉着子而来，便藉着那\Quote{在一切内工作万有}的圣神得以完成，他们为何反对自己的生命呢？为什么他们使自己疏离于\Quote{那些得救的人}所属的希望？为什么他们断绝自己与天主结合的关系？因为除非圣神在我们内运作我们与祂的结合，谁能与主结合呢？为什么他们与我们争论事奉和朝拜的多少？为什么他们以讽刺的贬义使用\Quote{朝拜}一词——这贬损了神圣而完全独立的存在者——，而设想他们渴望自己的救恩呢？我们要对他们说：你们的恳求是你们这些祈求者的益处，而非对那赐予者的尊崇。那么，为什么你们接近你们的恩主，仿佛你们有可给与之物？或者更准确地说，为什么你们全然拒绝称那赐予你们幸福者为恩主，轻视那赐生命者，却紧抓住生命呢？为什么你们寻求祂的圣化，却误解了圣化恩宠的分施者；至于那些祝福的给予，为什么你们不否认祂有能力，却认为祂不配被祈求施予，且未能考虑这一点：即给予某些祝福比起被请求给予是何等更大的事？祈求并不明确见证被祈求者的伟大；因为可能那没有东西可给的人也会被祈求，因为祈求只取决于祈求者的意愿。但实际施予某些祝福的人，由此给出了他内在之权能的无可置疑的证据。那么，既然你们为祂内那更大的事——我指赐予一切道德美善的能力——作证，为什么你们竟否认祂的祈求权，仿佛那是重要之事？虽然，如我们所说，祂的祈求常发生在那些自身无权能者身上，由于他们的崇拜者的迷惑。例如，迷信的奴仆向偶像祈求他们愿望的对象；但在此偶像的例子中，祈求并不赋予任何光荣；人们给予它们那种关注，只是由于错觉的期望——以为会得到他们所祈求的某物，因此他们不停地祈求。但是你们，既深信圣神是何种何等伟大之事的赐予者，竟忽视向祂祈求它们，而避难于那命令你们\Quote{只朝拜天主，且只事奉祂}的法律吗？那么，请告诉我，当你们将祂与祂自己的独生子和祂自己的圣神的内在结合割裂时，你们怎能单独朝拜祂？这种朝拜纯粹是犹太教式的。\par

但你会说：\Quote{当我想到圣父时，我已经把圣子（唯独）也包括在那个概念中了。}请告诉我；当你把握了圣子的概念时，难道你不也同时承认了圣神的概念吗？因为除非藉着圣神，你怎能宣认圣子呢？那么，圣神在什么时候是与圣子分离的，以致当圣父受到钦崇时，对圣神的钦崇不随着对圣子的钦崇一起包括在内？至于他们的钦崇本身，他们究竟将其看作什么？他们把它当作某种极致的荣耀献给万有的天主，有时也传递到独生子那里；但对圣神，他们却认为不配享有这样的特权。按人类通常的说法，下级在问候上级时俯伏于地的自谦姿态就称为钦崇。因此，先祖雅各伯在迎接兄弟并平息其愤怒时，似乎就用这种姿势来表示自己的卑微；因为圣经说\Quote{他向地下拜}，\Quote{三次}；若瑟的兄弟们，在他们尚未认识他、而他也在他们面前装作不认识他们时，因他地位的崇高，以这种钦崇来尊敬他的主权；甚至伟大的亚巴郎也\Quote{向赫特人下拜}，他是那地居民中的外方客旅，我认为，这个举动表明那些本地居民比寄居者强大得多。在古代记录中，以及在当今世界眼前的例子中，都可以说到许多这样的行为。\par

那么，他们所指的钦崇也是这个意思吗？认为用先祖甚至向客纳罕人表示敬意的礼来尊敬圣神是错误的，这岂非荒谬？或者他们认为他们的\Quote{钦崇}与此不同，似乎一种适合人，另一种适合至高者？但那么，他们为何在圣神身上完全省略了钦崇，连人在情况下所给予的钦崇也不赐予祂？他们想象专属于天主的哪种钦崇？是语言还是动作？但难道这些尊荣的标记不也为人类所共享吗？在人而言，也有语言和动作。那么，对于稍有反省能力的人，岂不明白人类并无任何配得上天主的恩赐可献；因为万福之源并不需要我们。而是我们人类将这些表达尊敬和仰慕的标记——我们在彼此之间采用，当我们要通过承认邻人的优越性来表示自己更为卑微时——转移到了对更高权能的侍奉上；我们从自己的所有中，将最珍贵的奉献给无价的本性。因此，既然人们因所求而接近君王和权贵时，不仅带去单纯的请求，还用尽一切可能的手段来引发他们的怜悯和恩惠，采用谦卑的声音、跪姿、抱膝、俯伏于地，并摆出各种可怜的表情来为请求说情，以唤醒同情——同样，那些认识真正权能者（祂掌管一切存在）的人，当他们为他们心中所想之事祈求时，有些人因今世的可怜境况而灵里卑微，有些人因永恒的奥秘盼望而思绪高昂，看到自己不知如何祈求，且他们的人性无法展示任何足以匹配那荣耀的伟大虔诚，便将用于人的礼仪搬到了对天主的侍奉上。这就是\Quote{钦崇}——即为了心中所求而连同祈求和自谦献上的东西。因此，达尼尔也向主屈膝，为俘虏的百姓祈求祂的爱；而\Quote{担负了我们的疾苦}、为我们代祷的那位，在福音中记载，因祂所取的人性，在祈祷时俯伏在地，以这种方式恳求，我想，这是命令我们在祈祷时声音不可大胆，而要采取可怜者的姿态；因为主\Quote{拒绝骄傲的人，却赐恩宠给谦卑的人}；在另一处祂说：\Quote{凡高举自己的，必被贬抑。}如果\Quote{钦崇}是一种祈祷的状态，或是为祈求的目标而提出的恳求，那么这些新规定有什么意图？这些人甚至不屑向赐予者祈求，不屑向统治者下跪，也不屑侍奉权能者。\par

\SourceRecord{Translated by William Moore and Henry Austin Wilson.；Nicene and Post-Nicene Fathers, Second Series；Vol. 5.；Edited by Philip Schaff and Henry Wace.；Buffalo, NY: Christian Literature Publishing Co.,；1893.；<http://www.newadvent.org/fathers/2903.htm>.}
